\section{A dedicated history selector: PA194B}
\label{sec:history-selector}

This construction gives a new physical decoder for the primary PA194 program.
It has 194 working controls and preserves the primary program's halting
behaviour on \emph{every} tape started in native control 74.  Its correctness
does not rely on bounded testing, on assumptions about reachable arithmetic
register values, or on assuming PA consistent.  It does rely on the inherited
calibration if that halting behaviour is identified with PA refutability.

\subsection{Exact construction}
Let $M$ denote the pinned primary PA194 table.  The relevant old rows, written
as $(\hbox{write},\hbox{move},\hbox{next})$, are
\[
\begin{array}{c|cc}
q & 0 & 1\\\hline
77 & (0,R,87)&(1,R,88)\\
87 & (0,R,301)&(1,R,21)\\
88 & (0,R,122)&(1,R,103)\\
301&(0,R,45)&(1,R,103)\\
122&(0,L,137)&(0,L,1)\\
137&(0,R,30)&(1,R,70)\\
30 &(0,R,45)&(1,R,46).
\end{array}
\]
The complete table has exactly one incoming edge to old control 88, namely
$77/1$, and exactly one incoming edge to old control 301, namely $87/0$.
Thus, immediately on entry to 88, the cell one position to the left contains
1; immediately on entry to 301, that cell contains 0.  The respective input
symbols have just been written, so the conclusion holds regardless of every
other tape cell.

Replace exactly four rows:
\[
\begin{array}{c|c|c}
\hbox{row}&M& M_B\\\hline
87/0 &(0,R,301)&(0,R,88)\\
88/0 &(0,R,122)&(0,L,301)\\
301/0&(0,R,45)&(0,R,122)\\
301/1&(1,R,103)&(1,R,301).
\end{array}
\]
Every other row, including every arithmetic-engine row, is unchanged.
In the new table, control 88 represents both old readers and control 301
is a dedicated history selector.  The reused numeral 301 is a label: its
new operational role is different from the old reader's role.

\subsection{Two-cell proof}
Take the entering head position to be zero.  Let $h$ be the symbol in cell
$-1$ and $b$ that in cell zero.  No assumption is made about any other cell.
For old control 88 use $h=1$, and for old control 301 use $h=0$.  In both cases
the new macro starts in control 88.  If $b=1$, both old readers have action
$(1,R,103)$, which is also the new $88/1$ action; both computations arrive
at control 103 at position 1 after one step.
If $b=0$, the following traces apply; each displayed number is a native
control and the parenthesized number is the head position:
\[
\begin{array}{c|l|c}
\hbox{old reader}&\hbox{new trace}&\hbox{duration}\\\hline
88&88(0),301(-1),301(0),122(1)&3\\
301&88(0),301(-1),122(0),137(-1),30(0),45(1)&5.
\end{array}
\]
All writes preserve the two symbols $h,b$.  Until the final configuration,
the head only visits cells $-1$ and zero.  Therefore these are proofs for
arbitrary complete tapes, not finite-window experiments that leave an
unexamined assumption about exterior cells.  The final head position,
control, and entire tape agree exactly with the corresponding old single
transition.  The new intermediate controls are all nonhalting.

\subsection{Whole-run theorem and internalization}
The start control 74 is distinct from both replaced readers.  For a run of
$M$ from control 74, align its configurations with a run of $M_B$: ordinary
transitions take one step, old 88/0 takes three steps, and old 301/0 takes
five steps.  The incoming-edge fact supplies the history bit whenever an
affected reader is entered.  Every aligned configuration has the same
head and tape, with old 301 represented by new 88.  Induction on the number
of old steps proves the alignment for every finite prefix.

Conversely, determinism and the four local cases partition the new run
into these finite nonhalting segments.  A halting new prefix must finish
at a segment boundary, and hence gives an old halting prefix.  This proves
\[
 \operatorname{Halt}(M_B,74,T,x)
 \quad\longleftrightarrow\quad
 \operatorname{Halt}(M,74,T,x)
\]
for every complete tape $T$ and initial head $x$.
If the old halting time is $t$, the aligned new halting time is
\[
 t+2a_t+4b_t\leq 5t,
\]
where $a_t$ and $b_t$ count visits to old rows $88/0$ and $301/0$ in its first
$t$ transitions.  Conversely the old halting time is at most the new one.
The step relation, the counting functions, and the segment expansion are
primitive recursive on finite configuration histories.  PA proves the
four local identities and their finite-run induction.  Consequently PA
proves the displayed equivalence for the blank tape as well.  Composing
with the inherited mathematical primary-PA194 calibration transfers
$\neg\operatorname{Halt}\leftrightarrow\operatorname{Con}(\mathrm{PA})$.
This argument is not a claim that the inherited calibration has acquired
an end-to-end Lean or iPRA proof.

\subsection{What changed and what was checked}
The change is a factoring of two recurring decoder readers through a
new history-sensitive selector.  Both successful paths and the extra
selector are exercised during the blank startup; the new helper is visited
55 times before the initial source frame.  This is an operational change
to recurring decoding, rather than a control renaming or an edit of an
unreached transition.  It gives no state saving: the dedicated selector
uses the slot released by joining the two old readers.

The tempting 193-control variant reuses arithmetic skip control 51 for
this selector.  Its $51/0$ rewrite changes the marker behaviour when the
following cell is blank.  Normal initialized register frames supply a 1
there, but blank initialization does not.  PA194B keeps control 51's two
original rows and avoids that exact obstruction.

The supplied executable checker reconstructs the dense table, verifies
all four two-cell identities, checks the complete incoming-edge lists,
and confirms unchanged arithmetic-engine rows.  It also runs all 4,780
symbolic branch outcomes for the 3,673 primitive source sites.  Its
complete blank-tape simulation reaches precisely the old source frame
at native control 84, head $-2$, with the same finite tape support and
register file, after 58,242 transitions; the old table requires 58,154.
The source and local arguments are the unbounded justification.  The
58,242-transition check establishes this particular initialization.

The dense table has states $0,\ldots,193$, initial state zero and separate
halt label 194, with exactly 388 transitions.  Its SHA-256 is
\begin{quote}\ttfamily\small
b2bcdb6c3816a1fbade699a79e431e4b530ec341aaf516c911920d43f9995397
\end{quote}
